\input{../style}

\title{{\Large Programmation sur Processeur Graphique -- GPGPU }\\
\vspace{-0.4em}
{\large TD 2b : outils de développement}\\
{\small Centrale Nantes}\\
\small P.-E. Hladik, \small{\texttt{pehladik@ec-nantes.fr}}\\
---\\
{\scriptsize  Version 1.0.0 (\today)}\\
}

\begin{document}

\maketitle


%%%%%%%%%%%%%%%%%%%%%%%%%
  \section{Utilisation de cuda-memcheck}
%%%%%%%%%%%%%%%%%%%%%%%%%

\begin{objectif}[]{}
 \begin{itemize}
	\item Utiliser cuda-memcheck
\end{itemize}
\end{objectif}

\begin{wtodo}[Utiliser {\tt memcheck} et ses options]{}

Commencer par lancer la commande {\tt sudo chmod a+rw /dev/nvhost-dbg-gpu} pour donner les droits à cuda-memcheck.

\begin{enumerate}
	\item  Compiler et exécuter le code {\tt deb1.cu}. Les résultats doivent être les chiffres de 0 à 9. Normalement vous devriez avoir une exécution correcte alors qu'il y a un débordement mémoire dans le kernel car ligne 7 si {\it i > N} on a {\tt a[i]} qui écrit à une adresse mémoire qui n'est pas allouée.
	
	Le C n'est pas capable de vérifier cela à l'exécution et ça ne pose pas de problème à l'exécution tant que l'on n'écrit à une adresse mémoire qui n'est pas accessible. Ce qui risque de poser problème avec un code plus important...
	\item  Exécuter {\tt cuda-memcheck ./a.out} qui vous signale des erreurs ({\tt ERROR SUMMARY: 6 errors}). En remontant dans la trace vous pouvez remarquer que {\tt memcheck} fait remonter des erreurs de type {\tt out of bounds} en indiquant dans quel thread l'erreur se produit (par exemple le thread 30 dans le bloc 1 : {\tt thread (30,0,0) in block (1,0,0)}). Par contre la trace n'est pas très lisible.
	\item Utiliser la commande {\tt cuda-memcheck -{}-show-backtrace no ./a.out} pour supprimer les traces arrières (et réduire les messages d'erreur).
	\item De même avec la commande {\tt cuda-memcheck -{}-report-api-errors no ./a.out} vous pouvez ne plus afficher les erreurs d'API (c'est-à-dire les appels à des fonctions qui ne sont extérieures à votre code). N'hésitez pas à combiner les deux options :
	
	{\tt cuda-memcheck -{}-show-backtrace no   -{}-report-api-errors no ./a.out}
	
	\item  Remarquez que {\tt memcheck} vous indique l'emplacement de l'erreur sous la forme d'une adresse dans le code compilé et non le code source (par exemple {\tt at 0x00000068}).  Pour faire le lien entre le code machine et le code source, il faut ajouter les options {\tt -Xcompiler -rdynamic -lineinfo} lors de la compilation avec {\tt nvcc} :
	
	 {\tt nvcc -Xcompiler -rdynamic -lineinfo deb1.cu}. 
	 
	 Après avoir recompilé avec ces options et en utilisant {\tt cuda-memcheck} vous devriez maintenant voir le numéro de ligne dans votre code source où a eu lieu l'erreur. Par exemple {\tt ...deb1.cu:7:foo(int*, int)} indique que l'erreur a eu lieu ligne 7 du fichier {\tt deb1.cu}.

	\item Modifiez la ligne qui pose problème en ajoutant un test {\tt if (i < N)} pour ne pas écrire dans une zone mémoire non allouée et vérifiez avec {\tt memcheck} que tout se passe bien.
\end{enumerate}

\end{wtodo}

%%%%%%%%%%%%%%%%%%%%%%%%%
  \section{Utilisation de gprof}
%%%%%%%%%%%%%%%%%%%%%%%%%

\begin{objectif}[]{}
 \begin{itemize}
	\item Utiliser gprog pour le profilage sur le CPU
\end{itemize}
\end{objectif}

\begin{wtodo}[Tester {\tt gprof}]{}
L'outil Gprof est un programme de profilage qui collecte et organise des statistiques sur d'un programme. Il ne permet de collecter que les informations sur le code qui s'exécute sur le CPU.
	
	\begin{enumerate}
		\item Compiler le code {\tt prof2.cu} avec l'option {\tt -pg} :
	
			{\tt nvcc -pg prof2.cu}
		
		La documentation sur {\tt -pg} indique que "-pg Generate extra code to write profile information suitable for the analysis program gprof. You must use this option when compiling the source files you want data about, and you must also use it when linking".
		\item Maintenant en exécutant le programme {\tt ./a.out}, un fichier {\tt gmon.out} contenant les mesures de performance sera généré.		
		\item  Utiliser {\tt gprof} sur le code {\tt prof2.cu} avec la commande  {\tt gprof -Q -b a.out}, les options {\tt -Q -b} permettent d'aller à l'essentiel, vous pouvez tester sans pour avoir tous les détails. Observer le temps nécessaire pour exécuter la fonction {\tt add\_v0}.
	\end{enumerate}

\end{wtodo}

\newpage

%%%%%%%%%%%%%%%%%%%%%%%%%
  \section{Utilisation de nvprof}
%%%%%%%%%%%%%%%%%%%%%%%%%

\begin{objectif}[]{}
 \begin{itemize}
	\item Utiliser nvprof
\end{itemize}
\end{objectif}

\begin{wtodo}[Tester {\tt nvprof}]{}

	\begin{enumerate}
		\item  Compiler le code {\tt prof2.cu} et collecter les informations de profilage en appelant {\tt nvprof} : \\{\tt sudo /usr/local/cuda/bin/nvprof ./a.out }
		\item  Comparer les performances de add\_v2 par rapport à add\_v1 ? A partir du code pouvez-vous expliquer ce que vous observez ?
		\item  Afficher les métriques disponibles : {\tt /usr/local/cuda/bin/nvprof -{}-query-metrics}
		\item  Afficher l'utilisation globale de la charge de calcul et de la mémoire : \\{\tt sudo /usr/local/cuda/bin/nvprof -m gld\_efficiency,gst\_efficiency ./a.out }
		\item Stocker la chronologie du profilage avec :\\ {\tt sudo /usr/local/cuda/bin/nvprof -o profile.timeline ./a.out}
		\item Enregistrez les mesures avec :\\ {\tt sudo /usr/local/cuda/bin/nvprof -o profile.metrics -{}-analysis-metrics ./a.out}
		%\item Regarder le fichier {\tt measure.h} et utiliser les macros pour mesurer le temps d'exécution de {\tt add\_v0}.
	\end{enumerate}

\end{wtodo}
\newpage

%%%%%%%%%%%%%%%%%%%%%%%%%
  \section{Utilisation de cuda-gdb (InfoIA: optionnel)}
%%%%%%%%%%%%%%%%%%%%%%%%%

\begin{objectif}[]{}
 \begin{itemize}
	\item Utiliser cuda-gdb
\end{itemize}
\end{objectif}

\begin{wtodo}[Test]{}

	\begin{enumerate}
		\item  Compiler et exécuter le code {\tt deb1.cu} avec les options {\tt -g} et {\tt -G} : {\tt cuda-gdb --args ./a.out}
		\item Exécutez quelques commandes dans cuda-gdb :
		\begin{verbatim}
			(cuda-gdb) b main                // set break point at main
			(cuda-gdb) r                     // run application
			(cuda-gdb) l                     // print line context 
			(cuda-gdb) b foo                 // break at kernel foo
			(cuda-gdb) c                     // continue
			(cuda-gdb) cuda thread           // print current thread
			(cuda-gdb) cuda thread 10        // switch to thread 10
			(cuda-gdb) cuda block            // print current block
			(cuda-gdb) cuda block 1          // switch to block 1
			(cuda-gdb) d                     // delete all break points
			(cuda-gdb) set cuda memcheck on  // turn on cuda memcheck
			(cuda-gdb) r                     // run from the beginning
		\end{verbatim}
		\item Retrouver le bug
	\end{enumerate}

\end{wtodo}



%%%%%%%%%%%%%%%%%%%%%%%%%%
%  \section{Bonus : utilisation de nvvp}
%%%%%%%%%%%%%%%%%%%%%%%%%%
%
%\begin{objectif}[]{}
% \begin{itemize}
%	\item Utiliser nvvp
%\end{itemize}
%\end{objectif}
%
%\begin{wtodo}[Test]{}
%	Installer nvvp \url{https://developer.nvidia.com/nvidia-visual-profiler}
%	\begin{enumerate}
%		\item Récupérez sur votre machine les fichiers {\tt profile.timeline} et {\tt profile.metrics}
%		\item Lancer nvvp puis
%			\begin{enumerate}
%				\item Choisissez File->Import... 
%				\item Sélectionnez Nvprof et cliquez sur Next,
%				\item Choississez Single process et cliquez sur Next,
%				\item Cliquez sur le 1er bouton Browse haut et sélectionnez le fichier {\tt profile.timeline},
%				\item Ajoutez les métriques à la ligne de temps en cliquant sur le 2ème bouton Browse et en sélectionnant le fichier {\tt profile.metrics},
%				\item Cliquez sur Finish
%			\end{enumerate}
%		\item Explorer la ligne de temps
%		\begin{itemize}
%			\item Contrôle + glisser de la souris dans la ligne de temps pour zoomer
%			\item Contrôle + glisser la souris dans la barre de mesure (en haut) pour mesurer le temps
%		\end{itemize}
%	\end{enumerate}
%
%\end{wtodo}

\end{document}